%% ch07_deployment.tex -- Intelligence Engineering, chapter 07 "Model
%% Deployment and Prediction Service".
%% Spine transferred from Huyen, Designing Machine Learning Systems, chapter
%% 7. The subchapters and the frame titles under them are the book's own
%% section and subsection headings, in the book's order.
%%
%% STATE: titles and TEMPLATE layout only. No body text. Every frame carries
%% the layout it is meant to be filled into, and a % TEMPLATE comment saying
%% why that layout and what belongs in each slot. Filling them is the next pass.
%%
%% Fragment of intelligence_engineering.tex. One chapter is one lecture and
%% gets exactly ONE \lecturepage, at the top. Inside it every \subsection
%% opens on the subchapter divider, which the master emits automatically via
%% \AtBeginSubsection; do not write those divider frames by hand.

\begin{refsection}

\lecturepage{VII}{Model Deployment and Prediction Service}{}
\section[Deployment]{Model Deployment and Prediction Service}

% ----------------------------------------------------------------------
\subsection{Machine Learning Deployment Myths}

% TEMPLATE: withmargin, the chapter's workhorse and the right opener for a
% myth: one claim the book contradicts plus one aside that does the
% contradicting. Main column states the myth and then the correction in the
% book's words. Margin takes the counts, the models behind a single ride
% sharing app and the company totals the book quotes.
\begin{frame}{Myth 1: You Only Deploy One or Two ML Models at a Time}
  \begin{withmargin}
    \begin{tdmain}
    \end{tdmain}
    \begin{tdmargin}
    \end{tdmargin}
  \end{withmargin}
  \slidesources{Huyen2022}
\end{frame}

% TEMPLATE: one plain box plus a takeaway, and deliberately not a second
% withmargin in a row. The heading already is the myth as a full sentence, so
% the box needs no header over it: it holds the myth and the reasons a
% deployed model decays. The takeaway carries the book's verdict on software
% rot, which only lands as a line of its own.
\begin{frame}{Myth 2: If We Don't Do Anything, Model Performance Remains the
    Same}
  \begin{tdblock}
  \end{tdblock}
  \begin{takeaway}
  \end{takeaway}
  \slidesources{Huyen2022}
\end{frame}

% TEMPLATE: withmargin. The heading is a claim about frequency, so the main
% column takes what forces an update and how often the book says it happens.
% Margin holds the named deployment cadences the book cites company by
% company, which are the evidence and read better stacked than in prose.
\begin{frame}{Myth 3: You Won't Need to Update Your Models as Much}
  \begin{withmargin}
    \begin{tdmain}
    \end{tdmain}
    \begin{tdmargin}
    \end{tdmargin}
  \end{withmargin}
  \slidesources{Huyen2022}
\end{frame}

% TEMPLATE: tddef, kept bare. This myth turns on one word the heading names,
% scale, and the book settles it by saying what scale means before arguing
% who has to worry about it. The term box is the slide's only slot: it takes
% the book's own measure of scale and the share of engineers it covers.
\begin{frame}{Myth 4: Most ML Engineers Don't Need to Worry About Scale}
  \begin{tddef}{Scale}
  \end{tddef}
  \slidesources{Huyen2022}
\end{frame}

% ----------------------------------------------------------------------
\subsection{Batch Prediction Versus Online Prediction}

% TEMPLATE: booktabs comparison table. The subchapter heading is the versus
% and the book already carries this contrast as a table, so it stays a table
% here rather than two boxes. The two column headers are the two regimes the
% heading names; the row labels are the book's own comparison dimensions and
% get filled together with the cells in the next pass.
\begin{frame}{From Batch Prediction to Online Prediction}
  \footnotesize
  \renewcommand{\arraystretch}{1.35}
  \begin{tabular}{@{}%
    >{\raggedright\arraybackslash}p{0.22\textwidth}%
    >{\raggedright\arraybackslash}p{0.35\textwidth}%
    >{\raggedright\arraybackslash}p{0.35\textwidth}@{}}
    \toprule
     & \textbf{Batch prediction} & \textbf{Online prediction}\\
    \midrule
     & & \\
     & & \\
     & & \\
     & & \\
    \bottomrule
  \end{tabular}
  \slidesources{Huyen2022}
\end{frame}

% TEMPLATE: side image hero, the first of this chapter's two. The heading asks
% an architecture question and the book answers it with a drawing of the two
% pipelines, so the strip takes that figure and the sidebody takes the steps
% in the book's order. The caption names the figure; the gray placeholder
% stays until the asset is drawn.
\begin{frame}{Unifying Batch Pipeline and Streaming Pipeline}
  \phimgside[0.33]{unifying_batch_pipeline_and_streaming_pipeline.png}{Batch pipeline and streaming pipeline}
  \begin{sidebody}
  \end{sidebody}
  \slidesources{Huyen2022}
\end{frame}

% ----------------------------------------------------------------------
\subsection{Model Compression}

% TEMPLATE: withmargin. The heading names one technique, so the main column
% takes what it replaces and the parameter arithmetic that makes it a
% compression method at all. Margin holds the architectures the book names as
% the worked cases, with their counts.
\begin{frame}{Low-Rank Factorization}
  \begin{withmargin}
    \begin{tdmain}
    \end{tdmain}
    \begin{tdmargin}
    \end{tdmargin}
  \end{withmargin}
  \slidesources{Huyen2022}
\end{frame}

% TEMPLATE: tddef. The heading is a named method the book introduces by
% definition, so the term box is the whole slide rather than a claim with an
% aside. The box takes the book's definition and the tradeoff it states
% immediately after it.
\begin{frame}{Knowledge Distillation}
  \begin{tddef}{Knowledge distillation}
  \end{tddef}
  \slidesources{Huyen2022}
\end{frame}

% TEMPLATE: two untitled boxes, levelled by the equal height group. The book
% gives pruning two distinct meanings, so each takes a box of its own instead
% of two bullets in one. No box headers here on purpose: the heading names
% only pruning, and the two senses arrive with the text in the next pass.
\begin{frame}{Pruning}
  \begin{columns}[T,onlytextwidth]
    \begin{column}{0.47\textwidth}
      \begin{tdblock}[equal height group=prune]
      \end{tdblock}
    \end{column}
    \begin{column}{0.47\textwidth}
      \begin{tdblock}[equal height group=prune]
      \end{tdblock}
    \end{column}
  \end{columns}
  \slidesources{Huyen2022}
\end{frame}

% TEMPLATE: code panel plus a takeaway. Quantization is a statement about how
% a number is represented, so the panel is the honest slot: it takes the bit
% widths and the ranges the book lists, which only make their point written
% out. The takeaway keeps the book's caveat about rounding and overflow.
\begin{frame}{Quantization}
  \begin{tdcode}
  \end{tdcode}
  \begin{takeaway}
  \end{takeaway}
  \slidesources{Huyen2022}
\end{frame}

% ----------------------------------------------------------------------
\subsection{ML on the Cloud and on the Edge}

% TEMPLATE: two block contrast, levelled by the equal height group. This frame
% opens the two deep headings under it, so it stays an overview and not a
% detail layout. The heading names two operations and each takes one box,
% compiling a model for a target and optimizing it for that target; the boxes
% carry the book's account of each step and the children carry the detail.
\begin{frame}{Compiling and Optimizing Models for Edge Devices}
  \begin{columns}[T,onlytextwidth]
    \begin{column}{0.47\textwidth}
      \begin{tdblockt}[equal height group=edge]{Compiling}
      \end{tdblockt}
    \end{column}
    \begin{column}{0.47\textwidth}
      \begin{tdblockt}[equal height group=edge]{Optimizing}
      \end{tdblockt}
    \end{column}
  \end{columns}
  \slidesources{Huyen2022}
\end{frame}

% TEMPLATE: side image hero, the second and last of this chapter's two. The
% book makes this heading concrete with a drawing of what an optimization
% does to a computation graph, so the strip takes that figure and the sidebody
% takes the named techniques under it. Caption repeats the heading as the
% figure's name; placeholder stays until the asset exists.
\begin{frame}{Model optimization}
  \phimgside[0.33]{model_optimization.png}{Model optimization}
  \begin{sidebody}
  \end{sidebody}
  \slidesources{Huyen2022}
\end{frame}

% TEMPLATE: withmargin. The heading is a turn of the argument rather than a
% set of peers, so the main column runs it as one claim: the search space, and
% what a learned model does inside it. Margin takes the named system the book
% uses as its case and the cost the book attaches to the search.
\begin{frame}{Using ML to optimize ML models}
  \begin{withmargin}
    \begin{tdmain}
    \end{tdmain}
    \begin{tdmargin}
    \end{tdmargin}
  \end{withmargin}
  \slidesources{Huyen2022}
\end{frame}

% TEMPLATE: withmargin with a captioned code panel in the main column. The
% book's browser material is named runtimes and libraries, so the panel is
% what the slide should show rather than a sentence about them. Margin takes
% the promise the book attaches to running in a browser and the limit it
% names in the same breath.
\begin{frame}{ML in Browsers}
  \begin{withmargin}
    \begin{tdmain}
      \begin{tdcodet}{Browser}
      \end{tdcodet}
    \end{tdmain}
    \begin{tdmargin}
    \end{tdmargin}
  \end{withmargin}
  \slidesources{Huyen2022}
\end{frame}

% ----------------------------------------------------------------------
\subsection{Summary}

% TEMPLATE: statement, kept bare. A summary is a closing line, not a slide of
% bullets, so it takes the loud centred form and nothing else. The main line
% is the book's own summary of what deployment costs; the substatement is the
% consequence it carries into the next chapter.
\begin{frame}{Summary}
  \begin{statement}
    \substatement{}
  \end{statement}
  \slidesources{Huyen2022}
\end{frame}

% ----------------------------------------------------------------------
% This chapter's own references. The refsection opened above scopes the
% citation numbering to this chapter, so chapter_pdfs/<this>.pdf resolves
% its own [n] markers instead of pointing at a list it does not contain.
\begin{frame}{References}
  \printbibliography[heading=none]
\end{frame}
\end{refsection}
